\documentclass[11pt]{article}
\usepackage[margin=1in]{geometry}
\usepackage{amsmath,amssymb}
\usepackage{hyperref}
\usepackage[numbers,sort&compress]{natbib}

\title{Topological Set Theoretic Estimation:\\
\large Follow-on Work of the Past Decade (2015--2026)}
\author{Project \texttt{thejeb}}
\date{\today}

\begin{document}
\maketitle

\begin{abstract}
A targeted survey of work in the decade since Carlson's 2012
topological-space recast of set theoretic estimation (STE)
\citep{carlson2012}.  Two findings dominate.  First, the
\emph{topological}, non-convex branch of STE is, as far as we can
determine, an almost single-author thread: Carlson and his direct
collaborators, with little independent uptake.  Second, no proof
assistant has mechanized the multi-set feasibility machinery that STE
sits on, which makes the present project's Lean development
first-of-kind.  Claims below are attributed to primary sources; nothing
here is asserted proven in our library unless it is machine-checked in
\texttt{Ste}.
\end{abstract}

\section{Carlson's own program (2014--2026)}
The dissertation grew into a sustained line of IEEE conference papers
with a rotating co-author cluster (early on R.\,E.\ Hiromoto and
R.\,B.\ Wells; later B.\ Ghosh, I.\,K.\ Dutta, M.\ Totaro,
S.\,R.\ Mikkilineni, B.\ Williams, M.\ Singh, and others).  The
consistent thesis: reduce every cipher to an equivalent \emph{substitution}
cipher, then attack it with STE, the AEP, and unicity distance.
\begin{itemize}
\item \citet{carlson2014breaking} --- journal version of the
  dissertation's known-ciphertext attack, with the original advisors.
\item \citet{ghosh2021isomorphic} --- product ciphers reduce to a single
  substitution cipher via a ``meta-alphabet'' construction.  This is the
  published counterpart of what we mechanize in
  \texttt{Ste.Carlson.Reduction} (Theorem~5.2, Corollary~5.3).
\item \citet{carlson2021keyspace} --- uses the isomorph count (the
  content of Carlson's Lemma~4.1, our \texttt{card\_consistent\_keys}) to
  shrink the effective brute-force key space.
\item \citet{carlson2022localentropy} --- the most substantive
  information-theoretic follow-on: Shannon's global unicity distance is
  an \emph{average}, not a lower bound, and a per-string \emph{local}
  unicity is defined.  A natural (deep) future mechanization target.
\item \citet{carlson2022equivalence} --- security implications of the
  product-to-substitution equivalence.
\item The mode-breaking series, culminating in
  \citet{hiromoto2025ctr}, ``Breaking the Counter (CTR) Mode.''
\item \citet{carlson2026entropy} --- a GoTech/University of Maryland
  technical report, the closest thing to a Carlson self-survey.  It
  restates STE (citing Combettes 1993 and Witsenhausen--Schweppe 1968),
  ties STE explicitly to the AEP and unicity, and frames STE as ``a
  precursor to the neural network.''  Notably it dates Carlson's
  \emph{topological}-space STE idea to \textbf{2006}, predating the
  dissertation.
\end{itemize}
We were unable to verify the patent portfolio attributed to Carlson in
secondary sources, and therefore cite none; this is an explicit gap.

\section{Independent uptake of topological STE}
We found \textbf{no evidence} of anyone outside Carlson's immediate
collaboration developing STE over topological (non-metric, non-convex)
spaces in this decade.  The classical Hilbert-space, convex-feasibility
STE lineage (Combettes 1993; see our companion survey) continues to be
cited widely, but those citations remain within the convex-optimization
tradition; none pivot to the topological generalization.  We record this
as a genuine gap and opportunity rather than padding it with tangential
citations.  (Absence of evidence from targeted search is not proof of
absence, but multiple queries returned nothing.)

\section{STE, the AEP, and unicity after 2012}
Post-2012 work tying STE to information theory --- property sets as AEP
typical sets, and unicity as feasibility collapse --- is substantial but
lives entirely inside the Carlson cluster, most explicitly in
\citet{carlson2022localentropy} and \citet{carlson2026entropy}.  We found
no independent work connecting STE specifically to the AEP, unicity, or
decryption in this window.

\section{Adjacent topological / feasibility work}
Being deliberately conservative: there is essentially no adjacent
literature that is genuinely ``STE in spirit, topological in substance''
outside Carlson's output.  Keyword-adjacent hits (e.g.\ constructions
using connected components of a topological space for encryption) are
cipher \emph{designs}, not estimation or feasibility methods, and we do
not claim an STE connection for them.  We deliberately exclude
persistent-homology / topological data analysis, a distinct mathematical
tradition with no feasibility link to STE.

\section{State of mechanization --- the opening}
This is the decisive input for the project's direction.
\begin{itemize}
\item \textbf{Lean / Mathlib.}  The single-set \emph{Hilbert projection
  theorem} is formalized (\texttt{Mathlib.Analysis.InnerProductSpace.\allowbreak
  Projection}): the existence and uniqueness of the nearest point in a
  nonempty complete convex set.  This is the base case of every
  projection method --- but \emph{no} iterative multi-set feasibility
  algorithm (alternating projections / POCS, Bregman cyclic projections,
  Douglas--Rachford) is in Mathlib.
\item \citet{li2024optlib} (Optlib) formalizes convex analysis,
  subgradients, proximal operators, and first-order convergence \emph{rates}
  on Hilbert spaces; \citet{cvxlean} (CvxLean) does convex problem
  \emph{modeling}.  Neither is feasibility- or projection-framed.
\item \textbf{Isabelle/HOL} (AFP) has strong Hilbert-space operator
  theory built for quantum verification, but nothing on projection-method
  convergence.  \textbf{Coq/Rocq} has convex polyhedra via the simplex
  method (finite-dimensional LP only).
\end{itemize}
Net: no proof assistant has POCS, Douglas--Rachford, Bauschke--Combettes
monotone-operator feasibility, or --- a fortiori --- topological STE.  The
single-set projection theorem is exactly the base case we would build on.
Our roadmap (POCS convergence over \texttt{Ste.Convex}, then the
topological generalization) is therefore unclaimed ground.

\section*{Method note}
This survey is exploratory and, where negative, honest about it: several
of the conclusions above are \emph{absence-of-evidence} findings from
targeted search, flagged as such.  As always, no mathematical claim is
treated as established for this project until it is machine-checked in the
\texttt{Ste} library.

\bibliographystyle{plainnat}
\bibliography{../../refs}

\end{document}
